\documentclass[11pt,a4paper]{article}
% Wspólny szablon dla kart supermocy (jedna karta = jedna strona A4)
% Użycie w karcie: \documentclass[11pt,a4paper]{article} % Wspólny szablon dla kart supermocy (jedna karta = jedna strona A4)
% Użycie w karcie: \documentclass[11pt,a4paper]{article} % Wspólny szablon dla kart supermocy (jedna karta = jedna strona A4)
% Użycie w karcie: \documentclass[11pt,a4paper]{article} \input{../preambula} \begin{document}

% Silnik tectonic (XeTeX): polskie znaki przez fontspec, nie przez inputenc/T1
\usepackage{fontspec}
\setmainfont{Arial}
\usepackage[polish]{babel}
\usepackage[a4paper,margin=1.5cm]{geometry}
\usepackage{amsmath,amssymb}
\usepackage{tikz}
\usetikzlibrary{arrows.meta,positioning,calc,shapes.geometric}
\usepackage[most]{tcolorbox}
\usepackage{enumitem}
\usepackage{booktabs}
\usepackage{xcolor}

\pagestyle{empty}
\setlength{\parindent}{0pt}
\setlength{\parskip}{0.4em}

\definecolor{cel}{HTML}{1F4E79}
\definecolor{regula}{HTML}{E8F1FA}
\definecolor{pulapka}{HTML}{FDECEA}
\definecolor{okres}{HTML}{C0392B}
\definecolor{skonczony}{HTML}{1E8449}

% Ramka na regułę
\newtcolorbox{regula}[1][]{
  colback=regula, colframe=cel, boxrule=0.8pt, arc=2pt,
  left=6pt, right=6pt, top=4pt, bottom=4pt, title=#1,
  fonttitle=\bfseries, coltitle=white
}

% Ramka na pułapkę z kartkówki
\newtcolorbox{pulapka}[1][]{
  colback=pulapka, colframe=okres, boxrule=0.8pt, arc=2pt,
  left=6pt, right=6pt, top=4pt, bottom=4pt, title=#1,
  fonttitle=\bfseries, coltitle=white
}

% Nagłówek karty: numer, tytuł, cel
\newcommand{\naglowekKarty}[3]{%
  {\LARGE\bfseries\color{cel} Karta #1: #2}\\[2pt]
  {\large Cel: #3}\par
  \vspace{0.3em}\noindent\rule{\linewidth}{1pt}\par
}

% Pole na obliczenia ucznia: ramka z delikatną kratką, argument = wysokość
\newcommand{\miejsce}[1]{%
  \par\vspace{0.2em}\noindent
  \begin{tikzpicture}
    \draw[step=0.5cm, gray!25, very thin] (0,0) grid (\textwidth,#1);
    \draw[gray!60, rounded corners=2pt] (0,0) rectangle (\textwidth,#1);
    \node[anchor=north west, font=\scriptsize, text=gray] at (0.05,#1) {obliczenia};
  \end{tikzpicture}\par}
 \begin{document}

% Silnik tectonic (XeTeX): polskie znaki przez fontspec, nie przez inputenc/T1
\usepackage{fontspec}
\setmainfont{Arial}
\usepackage[polish]{babel}
\usepackage[a4paper,margin=1.5cm]{geometry}
\usepackage{amsmath,amssymb}
\usepackage{tikz}
\usetikzlibrary{arrows.meta,positioning,calc,shapes.geometric}
\usepackage[most]{tcolorbox}
\usepackage{enumitem}
\usepackage{booktabs}
\usepackage{xcolor}

\pagestyle{empty}
\setlength{\parindent}{0pt}
\setlength{\parskip}{0.4em}

\definecolor{cel}{HTML}{1F4E79}
\definecolor{regula}{HTML}{E8F1FA}
\definecolor{pulapka}{HTML}{FDECEA}
\definecolor{okres}{HTML}{C0392B}
\definecolor{skonczony}{HTML}{1E8449}

% Ramka na regułę
\newtcolorbox{regula}[1][]{
  colback=regula, colframe=cel, boxrule=0.8pt, arc=2pt,
  left=6pt, right=6pt, top=4pt, bottom=4pt, title=#1,
  fonttitle=\bfseries, coltitle=white
}

% Ramka na pułapkę z kartkówki
\newtcolorbox{pulapka}[1][]{
  colback=pulapka, colframe=okres, boxrule=0.8pt, arc=2pt,
  left=6pt, right=6pt, top=4pt, bottom=4pt, title=#1,
  fonttitle=\bfseries, coltitle=white
}

% Nagłówek karty: numer, tytuł, cel
\newcommand{\naglowekKarty}[3]{%
  {\LARGE\bfseries\color{cel} Karta #1: #2}\\[2pt]
  {\large Cel: #3}\par
  \vspace{0.3em}\noindent\rule{\linewidth}{1pt}\par
}

% Pole na obliczenia ucznia: ramka z delikatną kratką, argument = wysokość
\newcommand{\miejsce}[1]{%
  \par\vspace{0.2em}\noindent
  \begin{tikzpicture}
    \draw[step=0.5cm, gray!25, very thin] (0,0) grid (\textwidth,#1);
    \draw[gray!60, rounded corners=2pt] (0,0) rectangle (\textwidth,#1);
    \node[anchor=north west, font=\scriptsize, text=gray] at (0.05,#1) {obliczenia};
  \end{tikzpicture}\par}
 \begin{document}

% Silnik tectonic (XeTeX): polskie znaki przez fontspec, nie przez inputenc/T1
\usepackage{fontspec}
\setmainfont{Arial}
\usepackage[polish]{babel}
\usepackage[a4paper,margin=1.5cm]{geometry}
\usepackage{amsmath,amssymb}
\usepackage{tikz}
\usetikzlibrary{arrows.meta,positioning,calc,shapes.geometric}
\usepackage[most]{tcolorbox}
\usepackage{enumitem}
\usepackage{booktabs}
\usepackage{xcolor}

\pagestyle{empty}
\setlength{\parindent}{0pt}
\setlength{\parskip}{0.4em}

\definecolor{cel}{HTML}{1F4E79}
\definecolor{regula}{HTML}{E8F1FA}
\definecolor{pulapka}{HTML}{FDECEA}
\definecolor{okres}{HTML}{C0392B}
\definecolor{skonczony}{HTML}{1E8449}

% Ramka na regułę
\newtcolorbox{regula}[1][]{
  colback=regula, colframe=cel, boxrule=0.8pt, arc=2pt,
  left=6pt, right=6pt, top=4pt, bottom=4pt, title=#1,
  fonttitle=\bfseries, coltitle=white
}

% Ramka na pułapkę z kartkówki
\newtcolorbox{pulapka}[1][]{
  colback=pulapka, colframe=okres, boxrule=0.8pt, arc=2pt,
  left=6pt, right=6pt, top=4pt, bottom=4pt, title=#1,
  fonttitle=\bfseries, coltitle=white
}

% Nagłówek karty: numer, tytuł, cel
\newcommand{\naglowekKarty}[3]{%
  {\LARGE\bfseries\color{cel} Karta #1: #2}\\[2pt]
  {\large Cel: #3}\par
  \vspace{0.3em}\noindent\rule{\linewidth}{1pt}\par
}

% Pole na obliczenia ucznia: ramka z delikatną kratką, argument = wysokość
\newcommand{\miejsce}[1]{%
  \par\vspace{0.2em}\noindent
  \begin{tikzpicture}
    \draw[step=0.5cm, gray!25, very thin] (0,0) grid (\textwidth,#1);
    \draw[gray!60, rounded corners=2pt] (0,0) rectangle (\textwidth,#1);
    \node[anchor=north west, font=\scriptsize, text=gray] at (0.05,#1) {obliczenia};
  \end{tikzpicture}\par}


\begin{document}

\naglowekKarty{1}{Odwrotność czy liczba przeciwna?}{odróżnić odwrotność od liczby przeciwnej i wyznaczać obie}

\vspace{0.3em}
\begin{regula}[REGUŁA – DWIE RÓŻNE OPERACJE]
\begin{tabular}{@{}p{3.6cm}p{5.2cm}p{6.6cm}@{}}
\textbf{Liczba przeciwna} & zmiana \textbf{znaku} & $-\frac{3}{7}$ jest przeciwna do $\frac{3}{7}$ \\[4pt]
\textbf{Odwrotność} & zamiana \textbf{licznika z mianownikiem} & $\frac{7}{3}$ jest odwrotnością $\frac{3}{7}$ \\
\end{tabular}
\end{regula}

\vspace{0.4em}
\textbf{Dlaczego odwrotność „robi fikołka”? Druga droga: sprawdzenie iloczynem.} Odwrotność to taka liczba, która razy dana liczba daje $1$:
$$\frac{3}{7}\cdot\frac{7}{3}=\frac{21}{21}=1 \qquad\qquad -\frac{3}{7}\cdot\left(-\frac{7}{3}\right)=1$$
Gdyby odwrotność miała inny znak niż liczba, iloczyn nie wyszedłby dodatni. Dlatego \emph{odwrotność zachowuje znak}, a \emph{przeciwna go zmienia}.

\begin{center}
\begin{tikzpicture}[>=Stealth, x=2.2cm, node distance=3.6cm,
  box/.style={draw, rounded corners=2pt, align=center, minimum height=0.9cm, minimum width=2.4cm, inner sep=3pt}]
  % Wiersz A: przeciwna = lustro względem zera
  \draw[->] (-3.4,0) -- (3.4,0) node[right] {$x$};
  \foreach \x in {-3,-2,-1,0,1,2,3} {
    \draw (\x,0.08) -- (\x,-0.08);
    \node[below, font=\small] at (\x,-0.1) {$\x$};
  }
  \fill[cel] (3/7,0) circle (2.2pt) node[below=5pt, font=\bfseries] {$\frac{3}{7}$};
  \fill[okres] (-3/7,0) circle (2.2pt) node[below=5pt, font=\bfseries, color=okres] {$-\frac{3}{7}$};
  \draw[->, thick, okres] (3/7,0.1) to[out=90, in=90, looseness=1.4] node[above, font=\small, color=okres] {lustro względem $0$} (-3/7,0.1);
  \node[font=\bfseries, anchor=west] at (-3.4,1.6) {A) przeciwna: ODBICIE przez zero};
\end{tikzpicture}
\end{center}

\vspace{-0.2em}
\begin{center}
\begin{tikzpicture}[>=Stealth, node distance=4.2cm,
  box/.style={draw, rounded corners=2pt, align=center, minimum height=0.9cm, minimum width=2.4cm, inner sep=3pt}]
  \node[font=\bfseries, anchor=west] at (-4.0,1.0) {B) odwrotność: FIKOŁEK (licznik $\leftrightarrow$ mianownik)};
  \node[box] (f1) at (-2.9,0) {$\dfrac{3}{7}$};
  \node[box, right=of f1] (f2) {$\dfrac{7}{3}=2\dfrac{1}{3}$};
  \draw[->, very thick, cel] (f1) -- node[above, font=\small\bfseries] {FIKOŁEK} (f2);
  \node[font=\small, below=6pt of f2, anchor=north] (chk) {sprawdzenie: $\frac{3}{7}\cdot\frac{7}{3}=1$};
\end{tikzpicture}
\end{center}

\vspace{0.2em}
\textbf{Przykłady rozstrzygnięte:}
\begin{center}
\renewcommand{\arraystretch}{1.4}
\begin{tabular}{@{}lllll@{}}
\toprule
Liczba & Przeciwna & Odwrotność & Co się zmieniło? \\
\midrule
$\dfrac{3}{7}$ & $-\dfrac{3}{7}$ & $\dfrac{7}{3}=2\dfrac{1}{3}$ & przeciwna: znak; odwrotność: fikołek \\
$-4$ & $4$ & $-\dfrac{1}{4}$ & $-4=-\dfrac{4}{1}$, więc fikołek daje $-\dfrac{1}{4}$ \\
$0{,}3$ & $-0{,}3$ & $\dfrac{10}{3}=3\dfrac{1}{3}$ & najpierw $0{,}3=\dfrac{3}{10}$ \\
$-2\dfrac{1}{3}$ & $2\dfrac{1}{3}$ & $-\dfrac{3}{7}$ & najpierw $-2\dfrac{1}{3}=-\dfrac{7}{3}$ \\
\bottomrule
\end{tabular}
\end{center}

\begin{pulapka}[PUŁAPKI Z KARTKÓWKI]
\begin{enumerate}[leftmargin=*, itemsep=2pt]
  \item \textbf{Liczba mieszana:} najpierw zamień na ułamek niewłaściwy. $1\dfrac{2}{5}=\dfrac{5\cdot 1+2}{5}=\dfrac{7}{5}$, więc odwrotność to $\dfrac{5}{7}$. NIGDY nie odwracaj $1\dfrac{2}{5}$ jako $1\dfrac{5}{2}$.
  \item \textbf{Ułamek dziesiętny:} najpierw zapisz go z kreską. $0{,}25=\dfrac{1}{4}$, więc odwrotność to $4$.
  \item \textbf{Liczba całkowita:} dopisz „niewidzialną $1$” na dole: $8=\dfrac{8}{1}$, więc odwrotność to $\dfrac{1}{8}$.
  \item \textbf{Zero} nie ma odwrotności, bo nie da się przez nie dzielić.
\end{enumerate}
\end{pulapka}

\end{document}
